\chapter{Code Structure and Development}
\label{ch:architecture}

This chapter is written for developers who want to understand, modify or
extend \pycmf{}.

\section{Repository layout}

\begin{lstlisting}
src/pycmf/
  __init__.py   version and top-level exports (stable methods only)
  core/         self-consistent driver, DIIS, occupation schemes,
                amplitudes, MPI helpers, experimental marker
  mol/          molecules: integrals, correlation potential, OBMP2,
                O2BMP2, OBDH, O2BDH, OO-DH, RO-OBMP2, NOCI
  downfold/     one-body downfolding, FCIDUMP export
  properties/   densities, dipole moments, <S^2>, hyperfine couplings,
                transition dipoles
  embed/        projection embedding, concentric localization (experimental)
  pbc/          k-point OBMP2, O2BMP2, quasiparticles, OBDF (experimental)
tests/          pytest suite
examples/       runnable scripts with reference output
benchmarks/     timing and gap scripts; inputs of version 0.2
docs/           theory, usage, validation, status, migration, this manual
\end{lstlisting}

\begin{longtable}{@{}p{0.30\textwidth} p{0.64\textwidth}@{}}
\caption{Modules of the package.}\\
\toprule
Module & Content \\
\midrule
\endfirsthead
\toprule
Module & Content \\
\midrule
\endhead
\bottomrule
\endfoot
\code{core/scf_loop.py} & self-consistent driver shared by all methods \\
\code{core/amplitudes.py} & first-order doubles amplitudes as spin blocks; \code{restrict_to_active} \\
\code{core/occupation.py} & occupation schemes \code{aufbau}, \code{mom}, \code{imom} \\
\code{core/mpi.py} & optional mpi4py helpers (\code{split}, \code{bcast}, \code{allreduce}) \\
\code{core/status.py} & \code{ExperimentalWarning} and \code{experimental(name)} \\
\code{mol/integrals.py} & integral backends \code{IncoreInts} (exact) and \code{DFInts} \\
\code{mol/bch.py} & correlation potential: first- and second-order BCH terms \\
\code{mol/obmp2.py} & \code{OBMP2Base}, \code{ROBMP2}, \code{UOBMP2}, factories \code{OBMP2}, \code{O2BMP2} \\
\code{mol/obdh.py} & \code{ROBDH}, \code{UOBDH}, factories \code{OBDH}, \code{O2BDH}, \code{FUNCTIONALS} \\
\code{mol/roobmp2.py} & \code{ROOBMP2} (experimental) \\
\code{mol/oodh.py} & \code{UOODH}, factory \code{OODH} (experimental): its own \code{build} returns the exact orbital gradient as the ov block of $\fbar$ \\
\code{mol/excited.py} & MOM-UHF excited determinants, channel selection \\
\code{mol/noci.py} & \code{NOCI}, \code{Determinant}, \code{make_basis} (experimental) \\
\code{downfold/obdf.py} & \code{OBDF} \\
\code{properties/rdm.py} & densities and dipole moments \\
\code{properties/spin.py} & $\langle S^2\rangle$ \\
\code{properties/hfcc.py} & isotropic hyperfine couplings \\
\code{properties/transition.py} & generalized Slater--Condon rules, transition dipoles (experimental) \\
\code{embed/partition.py} & SPADE partition, AO indices of atoms \\
\code{embed/cl.py} & concentric localization \\
\code{embed/projection.py} & \code{EmbeddedOBMP2}, \code{EmbeddedOBDH} \\
\code{pbc/kobmp2.py} & \code{KOBMP2}, \code{KO2BMP2}, \code{remove_linear_dependency_} \\
\code{pbc/kintegrals.py} & three-index GDF tensors in memory or HDF5 (\code{KDFInts}) \\
\code{pbc/kbch.py} & correlation potential with $k$-points \\
\code{pbc/ksymm.py} & irreducible $k$-points, expansion, symmetrization (\code{KSymmetry}) \\
\code{pbc/kipea.py} & static IP/EA corrections \\
\code{pbc/kqp.py} & second-order self-energy and quasiparticle energies \\
\code{pbc/kpeom.py} & partitioned EOM-MP2 \\
\code{pbc/kobdf.py} & \code{KOBDF} \\
\code{pbc/kuobmp2.py} & \code{KUOBMP2} (code of version 0.2) \\
\end{longtable}

\section{One self-consistent driver}

\code{core.scf_loop.kernel(mp, mo_energy, mo_coeff, nocc)} iterates over a
list of \emph{channels}. A channel is a spin channel for molecules (one for a
restricted, two for an unrestricted reference) and an irreducible $k$-point
for periodic systems. In each cycle the driver

\begin{enumerate}
  \item calls \code{mp.build(mo_energy, mo_coeff, nocc)}, which returns an
        \code{Intermediates} object with at least \code{e_tot},
        \code{fock_ao}, \code{fock_mo} and \code{dm_chan};
  \item prints the energy, $\Delta E$ and $\|\bar f_{ia}\|$ and checks
        convergence (both $|\Delta E|<$ \code{conv_tol} and
        $\|\bar f_{ia}\|<$ \code{conv_tol_grad});
  \item applies DIIS to the AO Fock matrices with the error vector $[F,D]$
        projected on a fixed orthonormal basis;
  \item diagonalises each Fock matrix in the span of the current orbitals,
        so that truncated or linearly dependent orbital spaces are preserved;
  \item occupies the new orbitals with \code{occupation.select}.
\end{enumerate}

It returns \code{(converged, e_tot, mo_energy, mo_coeff, intermediates)}.
The overlap matrix is taken from \code{mp.get_ovlp()} if the method defines
it (one per $k$-point) and from the mean field otherwise.

\section{The build step of molecular OBMP2}

\code{OBMP2Base.build} works in this order:

\begin{enumerate}
  \item HF Fock matrix from the current density (\code{get_hf()} returns the
        input mean field, or its HF counterpart for a Kohn--Sham object);
  \item integrals through \code{ao2mo()}: \code{IncoreInts} or
        \code{DFInts}, depending on \code{with_df};
  \item amplitude blocks: \code{core.amplitudes.make_blocks}, which applies
        \code{c_os}, \code{c_ss} and \code{level_shift};
  \item correlation potential and constant:
        \code{mol.bch.correlation_potential}, controlled by \code{bch_order};
  \item \code{get_reference(dm, h1e, vhf, fock_hf, e_hf)}, which returns the
        reference Fock matrix, the reference energy and the weight of $v$.
\end{enumerate}

The correlated Fock matrix is then the reference Fock matrix plus the weighted
potential. \code{ROBMP2} and \code{UOBMP2} differ only in \code{_occ} (2 or
1 electrons per occupied orbital of a channel) and in \code{_pack} and
\code{_unpack}, which convert between the arrays of \pyscf{} and the lists of
channels.

\section{Variants as hooks}

Variants of the method are hooks on this structure, not new modules:

\begin{itemize}
  \item \code{mol/obdh.py}: a mixin overrides \code{get_reference} so that
        it returns the Kohn--Sham Fock matrix and energy and the weight
        $a_c$.
  \item \code{mol/roobmp2.py}: changes the orbital update.
  \item \code{embed/projection.py}: adds the embedding potential in
        \code{get_reference} and computes the subsystem energy.
  \item \code{downfold/obdf.py}: calls \code{bch.correlation_potential}
        twice, with all amplitudes and with \code{restrict_to_active}.
  \item \code{pbc/kobmp2.py}: provides its own \code{build} and
        \code{get_ovlp} and reuses the driver with $k$-points as channels.
\end{itemize}

\section{Adding a method}

\begin{enumerate}
  \item Put the module at its final place (\code{mol/}, \code{pbc/} or
        \code{core/}) and call \code{pycmf.core.status.experimental(name)}
        in \code{__init__}. Do not export it in \code{pycmf/__init__.py}.
  \item Reuse \code{core.scf_loop.kernel} by providing
        \code{build(mo_energy, mo_coeff, nocc)}; derive from an existing
        class and override \code{get_reference} where possible.
  \item Add three kinds of tests: an exact limit (for example ``first cycle
        from HF orbitals equals \pyscf{} MP2'', ``all orbitals active equals
        FCI'' or ``$a_c=0$ is hybrid DFT''), a regression value, and input
        checks.
  \item When the method has been validated against published or independent
        results, remove the experimental marker, export it, and update
        \code{docs/status.md}.
\end{enumerate}

\section{Coding conventions}

\begin{itemize}
  \item \textbf{One method per module.} DF, DIIS, MOM, spin scaling and BCH
        order are attributes. Never add modules named \code{*_slow},
        \code{*_fast}, \code{*_mom} or \code{*_new}.
  \item \textbf{Placement.} Molecular code in \code{mol/}, periodic code in
        \code{pbc/}, shared code in \code{core/}. Reference implementations
        written as explicit loops belong in \code{tests/}.
  \item \textbf{\pyscf{} style.} Classes derive from
        \code{lib.StreamObject}; options are attributes listed in
        \code{_keys}; \code{kernel()} returns the main result; output goes
        through \code{pyscf.lib.logger}, never \code{print}.
  \item \textbf{Never modify the mean-field object} passed by the user. A test
        checks this.
  \item \textbf{Lint.} \code{ruff check src/pycmf tests examples} must pass
        (line length 110, single quotes). \code{src/pycmf/pbc} and
        \code{benchmarks} are excluded from ruff. The pre-commit hook runs
        \code{ruff format}.
\end{itemize}

\section{Development workflow}

\begin{lstlisting}[style=shell]
$ python3 -m venv .venv
$ env -u PYTHONPATH .venv/bin/pip install -e ".[dev]"
$ env -u PYTHONPATH .venv/bin/python -m pytest -m "not slow"   # before merging
$ env -u PYTHONPATH .venv/bin/python -m pytest                 # with slow tests
$ env -u PYTHONPATH .venv/bin/python -m pytest tests/test_obmp2.py::test_name
$ env -u PYTHONPATH .venv/bin/ruff check src/pycmf tests examples
$ cd examples && env -u PYTHONPATH ../.venv/bin/python check_examples.py --all
\end{lstlisting}

\code{pytest.ini} sets \code{pythonpath = src}, so the tests also run without
installing the package. The continuous integration
(\code{.github/workflows/ci.yml}) runs ruff and the fast tests on Python 3.9
and 3.12. Development happens on short-lived \code{feature/<name>} branches
that are rebased on the main branch; releases are tags.

\subsection*{Keeping the documentation consistent}

If a change alters behaviour, numbers or test counts, update the matching
document: \code{docs/theory.md} (equations), \code{docs/usage.md} (API),
\code{docs/status.md} (maturity and known issues),
\code{docs/validation.md} (what each test checks), \code{docs/kpoints.md}
($k$-point algorithm and timings), \code{docs/migration.md} (names of
version 0.2), and this manual (\code{docs/manual}).

\subsection*{Building this manual}

\begin{lstlisting}[style=shell]
$ cd docs/manual
$ latexmk -pdf pycmf_manual.tex       # or: pdflatex pycmf_manual (twice)
\end{lstlisting}

The manual uses only standard \LaTeX{} packages (\code{lmodern},
\code{booktabs}, \code{tabularx}, \code{longtable}, \code{listings},
\code{enumitem}, \code{fancyhdr}, \code{titlesec}, \code{hyperref}).
